%  Academic Article Template
%  Compile with: pdflatex -> bibtex -> pdflatex -> pdflatex

\documentclass[12pt, a4paper]{article}

% ── Packages ─────────────────────────────────────────────────
\usepackage[margin=1in]{geometry}
\usepackage{setspace}
\usepackage{titling}
\usepackage{abstract}
\usepackage{booktabs}
\usepackage{tabularx}
\usepackage{longtable}
\usepackage{array}
\usepackage{multirow}
\usepackage{amsmath, amssymb, amsthm}
\usepackage{graphicx}
\usepackage{float}
\usepackage{caption}
\usepackage{subcaption}
\usepackage{hyperref}
\usepackage[round, sort, authoryear]{natbib}
\usepackage{xcolor}
\usepackage{lmodern}
\usepackage[T1]{fontenc}
\usepackage[utf8]{inputenc}
\usepackage{microtype}
\usepackage{parskip}

% ── Hyperlink styling ─────────────────────────────────────────
\hypersetup{
    colorlinks   = true,
    linkcolor    = black,
    citecolor    = blue!70!black,
    urlcolor     = blue!70!black,
    pdftitle     = {Paper Title Here},
    pdfauthor    = {Author Name},
}

% ── Spacing ───────────────────────────────────────────────────
\doublespacing

% ── Theorem environments (optional) ──────────────────────────
\newtheorem{theorem}{Theorem}
\newtheorem{proposition}{Proposition}
\newtheorem{lemma}{Lemma}
\newtheorem{corollary}{Corollary}
\theoremstyle{definition}
\newtheorem{assumption}{Assumption}
\newtheorem{definition}{Definition}
\theoremstyle{remark}
\newtheorem{remark}{Remark}

% ── Custom column type for literature review table ────────────
\newcolumntype{P}[1]{>{\raggedright\arraybackslash}p{#1}}

%% ============================================================
%%  TITLE PAGE INFORMATION
%% ============================================================

\title{%
    \vspace{2cm}
    {\LARGE \textbf{Your Paper Title Goes Here}} \\[1.5em]
    {\large A Subtitle if Needed}
}
\author{%
    \textbf{Author Name} \\[0.4em]
    \small Department of Economics, University Name \\
    \small \href{mailto:author@university.edu}{author@university.edu}
}
\date{\small \today}

%% ============================================================
\begin{document}

% ── Title page ───────────────────────────────────────────────
\begin{titlepage}
    \maketitle
    \thispagestyle{empty}

    \vspace{2em}

    % ── Abstract ─────────────────────────────────────────────
    \begin{abstract}
        \noindent
        Replace this text with your abstract (150–250 words). The abstract should
        briefly state the purpose of the research, the principal results, and major
        conclusions. It should be self-contained and citation-free if possible.
    \end{abstract}

    \vspace{1em}

    % ── Keywords ─────────────────────────────────────────────
    \noindent\textbf{Keywords:} keyword one, keyword two, keyword three, keyword four

    \vspace{0.8em}

    % ── JEL Codes ────────────────────────────────────────────
    \noindent\textbf{JEL Classification:} C23, E44, G12   % Replace with your codes

    \vspace{2em}

    % ── Disclaimer ───────────────────────────────────────────
    \begin{small}
        \noindent\textbf{Disclaimer:} The views expressed in this paper are solely those
        of the author(s) and do not necessarily represent the views of any institution,
        employer, or funding body. The author(s) declare no conflicts of interest. Any
        errors or omissions are the sole responsibility of the author(s).
    \end{small}

\end{titlepage}

% ── Restart page numbering after title page ───────────────────
\setcounter{page}{1}
\pagenumbering{arabic}

%% ============================================================
\section{Introduction}
\label{sec:introduction}
%% ============================================================

Introduce the research question, motivation, and contribution of the paper here.
State why the topic is important and what gap in the literature this paper addresses.

Provide a brief road-map of the paper: Section~\ref{sec:literature} reviews the
relevant literature; Section~\ref{sec:theory} develops the theoretical framework;
Section~\ref{sec:data} describes the data; Section~\ref{sec:empirical} details the
empirical strategy; Section~\ref{sec:results} presents the results; and
Section~\ref{sec:conclusion} concludes.

%% ============================================================
\section{Literature Review}
\label{sec:literature}
%% ============================================================

Provide a narrative discussion of the related literature here, followed by the
summary table below.

\vspace{1em}

% ── Literature Review Table ───────────────────────────────────
\begin{table}[H]
    \centering
    \caption{Summary of Related Literature}
    \label{tab:literature}
    \small
    \begin{tabularx}{\textwidth}{P{2.8cm} P{1.4cm} P{2.6cm} P{2.6cm} P{3.8cm}}
        \toprule
        \textbf{Study} & \textbf{Year} & \textbf{Data / Sample} & \textbf{Method} & \textbf{Main Finding} \\
        \midrule
        Author A \& Author B  & 20XX & Country/period & OLS / IV        & Brief finding here. \\[4pt]
        Author C              & 20XX & Country/period & Panel FE        & Brief finding here. \\[4pt]
        Author D et al.       & 20XX & Country/period & Diff-in-Diff    & Brief finding here. \\[4pt]
        Author E \& Author F  & 20XX & Country/period & RDD             & Brief finding here. \\[4pt]
        Author G              & 20XX & Country/period & ML / Structural & Brief finding here. \\
        \bottomrule
    \end{tabularx}
    \begin{minipage}{\textwidth}
        \vspace{3pt}
        \footnotesize\textit{Note:} OLS = Ordinary Least Squares; IV = Instrumental Variables;
        FE = Fixed Effects; Diff-in-Diff = Difference-in-Differences; RDD = Regression
        Discontinuity Design; ML = Machine Learning.
    \end{minipage}
\end{table}

%% ============================================================
\section{Theoretical Framework}
\label{sec:theory}
%% ============================================================

Present the theoretical model or conceptual framework here. Use numbered equations
and refer to them in the text.

\subsection{Model Setup}

Let $i = 1, \dots, N$ index agents and $t = 1, \dots, T$ index time periods.
Agent $i$ maximises the objective function:

\begin{equation}
    \max_{\{x_{it}\}} \sum_{t=0}^{T} \beta^{t} u(x_{it}, z_{it})
    \label{eq:objective}
\end{equation}

\noindent subject to the constraint:

\begin{equation}
    c_{it} + s_{it} = w_{it} + (1 + r_{t}) s_{it-1}
    \label{eq:budget}
\end{equation}

\subsection{Equilibrium and Predictions}

State the key theoretical predictions / testable hypotheses that follow from the
model. These will be linked to the empirical tests in Section~\ref{sec:empirical}.

\begin{proposition}
    Under Assumptions 1–3, the equilibrium outcome satisfies $\partial y^{*} / \partial x > 0$.
\end{proposition}

\begin{proof}
    Proof sketch here.
\end{proof}

%% ============================================================
\section{Data}
\label{sec:data}
%% ============================================================

Describe the data sources, sample construction, and key variables.

\subsection{Data Sources}

The primary data come from [Source]. The sample covers [time period] and [geographic
scope]. After removing observations with missing values on key variables, the final
sample comprises [N] observations.

\subsection{Variable Definitions and Descriptive Statistics}

\begin{table}[H]
    \centering
    \caption{Descriptive Statistics}
    \label{tab:descriptives}
    \begin{tabular}{lrrrrrr}
        \toprule
        \textbf{Variable} & \textbf{N} & \textbf{Mean} & \textbf{Std.\ Dev.}
                          & \textbf{Min}  & \textbf{Median} & \textbf{Max} \\
        \midrule
        Dependent Var.  & 000 & 0.000 & 0.000 & 0.000 & 0.000 & 0.000 \\
        Main Regressor  & 000 & 0.000 & 0.000 & 0.000 & 0.000 & 0.000 \\
        Control 1       & 000 & 0.000 & 0.000 & 0.000 & 0.000 & 0.000 \\
        Control 2       & 000 & 0.000 & 0.000 & 0.000 & 0.000 & 0.000 \\
        Control 3       & 000 & 0.000 & 0.000 & 0.000 & 0.000 & 0.000 \\
        \bottomrule
    \end{tabular}
    \begin{minipage}{\textwidth}
        \vspace{3pt}
        \footnotesize\textit{Note:} Describe variable definitions and units here.
    \end{minipage}
\end{table}

%% ============================================================
\section{Empirical Strategy}
\label{sec:empirical}
%% ============================================================

Describe the identification strategy and econometric specification.

\subsection{Baseline Specification}

The baseline estimating equation is:

\begin{equation}
    y_{it} = \alpha + \beta \, x_{it} + \gamma' \mathbf{Z}_{it} + \mu_{i} + \lambda_{t} + \varepsilon_{it}
    \label{eq:baseline}
\end{equation}

\noindent where $y_{it}$ is the outcome variable, $x_{it}$ is the main explanatory
variable of interest, $\mathbf{Z}_{it}$ is a vector of controls, $\mu_{i}$ denotes
unit fixed effects, $\lambda_{t}$ denotes time fixed effects, and $\varepsilon_{it}$
is the error term. Standard errors are clustered at the [level] level.

\subsection{Identification}

Discuss potential endogeneity concerns and how the empirical strategy addresses them
(e.g., instrumental variables, difference-in-differences, RDD, etc.).

%% ============================================================
\section{Results}
\label{sec:results}
%% ============================================================

Present and discuss the main empirical findings.

\subsection{Baseline Results}

\begin{table}[H]
    \centering
    \caption{Baseline Regression Results}
    \label{tab:results_main}
    \begin{tabular}{lcccc}
        \toprule
        & \multicolumn{4}{c}{\textit{Dependent variable: } $y_{it}$} \\
        \cmidrule(lr){2-5}
        & (1) & (2) & (3) & (4) \\
        \midrule
        $x_{it}$       & 0.000***  & 0.000***  & 0.000**   & 0.000*    \\
                        & (0.000)   & (0.000)   & (0.000)   & (0.000)   \\[4pt]
        Control 1       &           & 0.000     & 0.000     & 0.000     \\
                        &           & (0.000)   & (0.000)   & (0.000)   \\[4pt]
        Control 2       &           &           & 0.000     & 0.000     \\
                        &           &           & (0.000)   & (0.000)   \\[4pt]
        Constant        & 0.000***  & 0.000***  & 0.000***  & 0.000***  \\
                        & (0.000)   & (0.000)   & (0.000)   & (0.000)   \\
        \midrule
        Unit FE         & No        & No        & Yes       & Yes       \\
        Time FE         & No        & No        & No        & Yes       \\
        Observations    & 000       & 000       & 000       & 000       \\
        $R^{2}$         & 0.000     & 0.000     & 0.000     & 0.000     \\
        \bottomrule
    \end{tabular}
    \begin{minipage}{\textwidth}
        \vspace{3pt}
        \footnotesize\textit{Note:} Standard errors clustered at the [level] level in
        parentheses. *** $p<0.01$, ** $p<0.05$, * $p<0.10$.
    \end{minipage}
\end{table}

\subsection{Robustness Checks}

Describe robustness checks, alternative specifications, placebo tests, etc.

%% ============================================================
\section{Conclusion}
\label{sec:conclusion}
%% ============================================================

Summarise the main findings and their implications. Discuss the contributions to
the literature and policy relevance. Acknowledge limitations and suggest directions
for future research.

%% ============================================================
%%  REFERENCES
%% ============================================================

\newpage
\singlespacing
\bibliographystyle{apalike}   % or: plainnat, apa, chicago, etc.
\bibliography{references}     % links to references.bib

%% ============================================================
%%  APPENDIX (optional)
%% ============================================================

\newpage
\appendix
\doublespacing

\section{Additional Tables and Figures}
\label{app:tables}

Place supplementary material here.

\section{Proofs}
\label{app:proofs}

Place extended proofs here.

\end{document}